% Problem 1: Structured report content (placeholder figures to be inserted)
\section{Autoencoder / MLP Analysis}

\subsection{Dataset and Preprocessing}
The dataset used, directory layout, and preprocessing steps are listed below:
\begin{itemize}
  \item Grayscale input conversion to single channel
  \item Rescaling / normalization: placeholder for mean/std (to be filled)
  \item Train / test split as provided in the assignment dataset
\end{itemize}

\subsection{Model Architecture and Decisions}
We implemented an MLP / autoencoder-style architecture. Key design decisions:
\begin{itemize}
  \item Hidden layer sizes and activation functions selected to balance capacity and training stability.
  \item Loss: cross-entropy for classification tasks or mean-squared error for reconstruction tasks, denoted here as $\mathcal{L}$.
\end{itemize}

\subsubsection{Mathematical Formulation}
For classification we use softmax cross-entropy. Given logits $z\in\mathbb{R}^C$ for $C$ classes, the predicted probabilities are
$$p_c = \frac{e^{z_c}}{\sum_{j=1}^C e^{z_j}},$$
and the loss for a single example with one-hot target $y$ is
$$\mathcal{L}_{CE} = -\sum_{c=1}^C y_c \log p_c.$$ 
For reconstruction tasks with predicted output $\hat{x}$ and target $x$, we use MSE:
$$\mathcal{L}_{MSE} = \|x - \hat{x}\|_2^2.$$

\subsection{Training Procedure}
Hyperparameters used (placeholders; fill with actual values from runs):
\begin{itemize}
  \item Optimizer: SGD/Adam
  \item Learning rate: 0.001 (example)
  \item Batch size: 128
  \item Epochs: 5
\end{itemize}
Training and evaluation routines measured wall-clock time using high-precision timers; reported metrics include training accuracy, test accuracy, training time, and test time.

\subsection{Results}
Insert accuracy and timing tables and plots here. Example table template:

\begin{table}[ht]
\centering
\caption{Problem 1: Summary of results}
\begin{tabular}{lrrr}
\hline
Model & Train Acc (\%) & Test Acc (\%) & Train Time (s) \\
\hline
Baseline MLP & -- & -- & -- \\
Variant A & -- & -- & -- \\
\hline
\end{tabular}
\end{table}

% Placeholder for figures
\begin{figure}[ht]
  \centering
  \fbox{\parbox[c][4cm][c]{0.7\textwidth}{\centering Placeholder: prob1\_train\_loss.png}}
  \caption{Training loss curve (Problem 1).}
  \label{fig:prob1_loss}
\end{figure}

\subsection{Discussion and Decisions}
Summarize observed behaviour, issues encountered (e.g., GPU compatibility), and rationale for chosen hyperparameters.

\subsection{Conclusion}
State final conclusions and next steps.

% End of Problem 1
